\chapter{Brain Abscess}

\section*{Definition and Overview}
A brain abscess is a localised collection of pus that forms within the substance of the brain, usually as the result of a bacterial, fungal, or parasitic infection. The process begins with an area of infected inflammation known as cerebritis, which over time matures into a walled-off pocket of necrotic material surrounded by a fibrous capsule. Because the skull is a rigid, fixed-volume container, even a relatively small abscess can become life-threatening by compressing adjacent brain tissue, obstructing the flow of cerebrospinal fluid, and raising the pressure inside the head. Brain abscesses are uncommon but constitute a true neurological emergency: they carry a real risk of permanent neurological injury or death, and every hour of delay in recognition and treatment worsens the outcome. With modern imaging, neurosurgical drainage, and prolonged antibiotics, however, most patients now survive, and many recover well.

\section*{Epidemiology}
Brain abscess is a rare condition, with reported incidences of roughly 0.4--0.9 cases per 100,000 population per year in developed countries. It can occur at any age but is most frequently seen in children, young adults, and older individuals with relevant risk factors. Males are affected approximately two to three times more often than females. The epidemiology has shifted over time in line with changes in predisposing conditions: in developed countries the classic sources -- chronic otitis media, paranasal sinusitis, and dental sepsis -- remain common, but an increasing proportion of cases now arises in immunocompromised patients and after neurosurgical procedures or penetrating trauma. In low- and middle-income countries, ear and sinus infections still account for the majority of cases, often presenting later and with more advanced disease. Mortality has fallen dramatically, from figures approaching 40\% in older series to roughly 5--15\% in contemporary series, reflecting better imaging, earlier diagnosis, and improved antimicrobial and surgical care.

\section*{Aetiology and Pathophysiology}
A brain abscess develops when organisms gain access to the brain parenchyma by one of three principal routes:

\begin{itemize}
    \item \textbf{Direct (contiguous) spread:} infection extends from an adjacent focus such as otitis media or mastoiditis (typically producing abscesses in the temporal lobe or cerebellum), frontal or ethmoid sinusitis (frontal lobe), or dental sepsis
    \item \textbf{Haematogenous spread:} organisms travel in the bloodstream from a distant source such as infective endocarditis, a lung abscess or bronchiectasis, or an infected skin wound -- this route predominates in intravenous drug users and commonly produces multiple abscesses at the grey--white matter junction
    \item \textbf{Direct inoculation:} a penetrating head injury, a compound skull fracture, or a neurosurgical procedure introduces organisms directly into the brain
\end{itemize}

The causative organisms vary with the route and the host. Streptococci (including the \textit{Streptococcus milleri} group), staphylococci, and anaerobic bacteria are the most common isolates, often in mixed culture. In immunocompromised patients, \textit{Toxoplasma gondii}, \textit{Nocardia}, fungi, and mycobacteria must be considered. Neonates are more often infected with Gram-negative organisms such as \textit{Citrobacter} and \textit{Proteus}. Pathologically, the lesion evolves through an early cerebritis phase, during which the inflammatory response produces oedema and tissue necrosis, to a late encapsulated phase in which a collagenous capsule forms around the liquefied centre. Surrounding vasogenic oedema, the mass effect of the pus, and the irritant effect of infection on adjacent structures all contribute to raised intracranial pressure, focal deficits, and seizures.

\section*{Clinical Features}
The presentation is typically subacute, evolving over days to a few weeks, and varies with the site, size, and number of lesions. Common features include:

\begin{itemize}
    \item \textbf{Headache:} the most frequent symptom, often dull and persistent, and typically progressive
    \item \textbf{Fever and malaise:} present in a majority but by no means all patients, and less reliable in the immunocompromised
    \item \textbf{Nausea and vomiting:} reflecting raised intracranial pressure
    \item \textbf{Focal neurological deficits:} such as hemiparesis, hemisensory loss, aphasia, or visual field defects, depending on the location of the abscess
    \item \textbf{Seizures:} either focal or generalised, occurring in roughly 25--35\% of patients
    \item \textbf{Cognitive and behavioural change:} confusion, drowsiness, disorientation, or personality change, which may dominate in older patients
    \item \textbf{Neck stiffness:} if the infection irritates the meninges, although this is less prominent than in frank meningitis
\end{itemize}

Because the signs are non-specific and the evolution gradual, brain abscess is often initially mistaken for stroke, tumour, or migraine. Cerebellar abscesses present particular difficulty, with ataxia, nystagmus, and incoordination rather than hemiparesis.

\section*{Differential Diagnosis}
The differential diagnosis includes other intracranial space-occupying and infective processes:

\begin{itemize}
    \item Brain tumour (primary glioma or metastasis) -- may be indistinguishable on imaging alone
    \item Meningitis or encephalitis -- more acute onset with prominent fever and meningism
    \item Subdural empyema -- a collection of pus over the brain surface, often of similar origin
    \item Extradural abscess -- pus outside the dura, usually post-operative or post-traumatic
    \item Ischaemic or haemorrhagic stroke, including haemorrhage into a tumour
    \item Cerebral toxoplasmosis, tuberculoma, or cryptococcoma in immunocompromised hosts
    \item Septic cerebral venous sinus thrombosis
    \item Neurocysticercosis in endemic regions or in people who have travelled there
\end{itemize}

\section*{When to Seek Medical Care}
Anyone with a persistent or worsening headache, especially with fever, vomiting, new neurological symptoms, or drowsiness, should be assessed urgently. People known to have risk factors such as sinusitis, middle ear disease, dental infection, a heart valve infection, or a weakened immune system should not dismiss such symptoms as routine.

\textbf{Call 000 (or local emergency services) for:}
\begin{itemize}
    \item A sudden, severe headache -- the worst headache you have ever had
    \item Any seizure or fit
    \item Sudden weakness, numbness, difficulty speaking, or confusion
    \item Drowsiness or difficulty waking a person
    \item Neck stiffness together with fever and headache
    \item Collapse or loss of consciousness
\end{itemize}

\section*{Diagnosis and Investigations}
\begin{itemize}
    \item \textbf{History:} careful enquiry about fevers, ear and sinus symptoms, dental problems, recent surgery or head injury, intravenous drug use, immunosuppression, and travel history
    \item \textbf{Examination:} a full neurological assessment including fundoscopy to look for papilloedema, plus examination of the ears, sinuses, teeth, and heart
    \item \textbf{Blood tests:} full blood count, inflammatory markers (CRP, ESR), blood cultures, and -- in immunocompromised patients -- HIV serology and other targeted tests
    \item \textbf{Imaging:} contrast-enhanced CT is the usual first test and typically shows a ring-enhancing lesion with surrounding oedema; MRI with diffusion-weighted imaging is more sensitive and can distinguish abscess from tumour and detect cerebritis early
    \item \textbf{Microbiological sampling:} culture of pus obtained by stereotactic aspiration or surgical excision is the gold standard for identifying the organism and guiding antimicrobial therapy; if the patient is stable, the lumbar puncture findings may support an alternative diagnosis, but lumbar puncture is generally avoided once a space-occupying lesion is suspected
    \item \textbf{Source identification:} imaging of the sinuses, mastoids, teeth, and chest, and echocardiography if endocarditis is suspected, to find and treat the origin of infection
\end{itemize}

\section*{Management}
Brain abscess is managed in hospital as an emergency, usually by a multidisciplinary team of neurosurgeons, neurologists, and infectious diseases physicians:

\begin{itemize}
    \item \textbf{Antibiotics:} prolonged intravenous therapy (commonly 6--8 weeks) started empirically to cover streptococci, staphylococci, and anaerobes, then narrowed according to culture results; specific regimens are required for fungal or parasitic causes
    \item \textbf{Surgical drainage:} stereotactic aspiration is the standard approach for most abscesses, allowing both decompression and diagnosis; excision is reserved for selected cases, and small, multiple, or inaccessible lesions may occasionally be managed with antibiotics alone under close radiological follow-up
    \item \textbf{Control of raised intracranial pressure:} corticosteroids such as dexamethasone are used judiciously to reduce surrounding oedema, and measures to reduce ICP are employed if the pressure is critically elevated
    \item \textbf{Antiepileptic drugs:} given to treat or, in some cases, to prevent seizures during the acute phase
    \item \textbf{Treatment of the source:} drainage of infected sinuses, mastoidectomy, dental extraction, or valve surgery as appropriate, to prevent recurrence
    \item \textbf{Supportive care:} fluids, analgesia, antipyretics, and repeated neurological observation, with serial imaging to confirm resolution
\end{itemize}

\section*{Complications}
\begin{itemize}
    \item Raised intracranial pressure leading to brain herniation -- the most dangerous early complication
    \item Rupture of the abscess into the ventricular system, causing severe ventriculitis and an abrupt clinical deterioration
    \item Hydrocephalus, particularly with cerebellar abscesses compressing the fourth ventricle
    \item Cerebral oedema and infarction of adjacent brain tissue
    \item Seizures, which may persist as epilepsy after recovery
    \item Permanent focal deficits such as hemiparesis, aphasia, or visual loss
    \item Septicaemia and systemic sepsis with multi-organ dysfunction
    \item Recurrence if the source of infection is not eradicated or the abscess is inadequately drained
\end{itemize}

\section*{Prognosis and Outcomes}
With prompt recognition and modern treatment, the mortality of brain abscess is now approximately 5--15\%, a dramatic improvement over historic figures. Full recovery is possible, but a substantial proportion of survivors are left with permanent sequelae -- most commonly epilepsy, focal weakness, cognitive impairment, or memory difficulties. Prognosis is worse when diagnosis is delayed, when the patient is immunocompromised or deeply comatose at presentation, when the abscess has ruptured into the ventricles, or when there are multiple large lesions. Functional outcome depends heavily on the site of the abscess and the quality of rehabilitation. Survivors require neurological follow-up, often including antiepileptic medication and structured rehabilitation with physiotherapy, occupational therapy, and neuropsychology.

\section*{Prevention and Self-Care}
\begin{itemize}
    \item Treat ear, sinus, and dental infections promptly rather than allowing them to smoulder
    \item Maintain good dental hygiene, including regular check-ups, and seek early care for dental pain or abscesses
    \item Wear protective headgear for activities with a risk of head injury, and seek urgent review of any deep, dirty, or penetrating head wound
    \item Follow advice about the risks of injecting drug use, and seek help if you use intravenous drugs
    \item Manage conditions that increase risk, such as diabetes, and keep immunosuppressive medicines under specialist review
    \item After treatment for an abscess, attend all follow-up scans and appointments and report any new headache, seizure, or neurological symptom immediately
    \item If you have structural heart disease or a prosthetic valve, follow your cardiologist's advice about preventing blood-borne infection
\end{itemize}

\section*{Sources and Further Reading}
The following organisations provide reliable, evidence-based information for patients, students, and clinicians.

\begin{itemize}
    \item NHS. \textit{Brain abscess: symptoms, diagnosis, and treatment overview.} https://www.nhs.uk/conditions/
    \item Mayo Clinic. \textit{Brain abscess: patient education information.} https://www.mayoclinic.org/diseases-conditions/
    \item MSD Manual Professional Version. \textit{Brain abscess: aetiology, presentation, and management.} https://www.msdmanuals.com/
    \item MedlinePlus. \textit{Brain abscess: consumer health summary.} https://medlineplus.gov/
    \item Centers for Disease Control and Prevention. \textit{Information on bacterial and fungal infections of the central nervous system.} https://www.cdc.gov/
    \item World Health Organization. \textit{Global guidance on neurological infections and antimicrobial resistance.} https://www.who.int/
\end{itemize}
